\chapter{กำลังและเครื่องกลอย่างง่าย}
\label{ch:ch08_power_simple_machines}
\index{กำลังและเครื่องกล}

\begin{conceptbox}[title=วัตถุประสงค์การเรียนรู้ประจำบท]
เมื่อสิ้นสุดการศึกษาบทเรียนนี้ นักศึกษาสามารถ
\begin{enumerate}
    \item เมื่อศึกษาบทนี้จบแล้ว ผู้เรียนควรมีความสามารถดังต่อไปนี้
    \item อธิบายความหมายของกำลังเชิงกล (mechanical power) และความสัมพันธ์ระหว่างกำลังกับงานและเวลา
    \item คำนวณกำลังของระบบต่างๆ โดยใช้สูตร P = W/t และ P = Fv พร้อมทั้งแปลงหน่วยระหว่างวัตต์ (W) กับแรงม้า (hp)
    \item อธิบายหลักการทำงานของเครื่องกลอย่างง่ายทั้ง 6 ชนิด ได้แก่ คาน รอก พื้นเอียง ล้อและเพลา สกรู และลิ่ม
    \item คำนวณค่าได้เปรียบเชิงกล (mechanical advantage, MA) ทั้งในเชิงทฤษฎีและเชิงปฏิบัติของเครื่องกลแต่ละชนิด
\end{enumerate}
\end{conceptbox}

\section{ทฤษฎีและหลักการพื้นฐาน}

\begin{bioappbox}
\textbf{กำลังชีวภาพและประสิทธิภาพเชิงกลของเครื่องจักรกลการเกษตร}

อัตราการทำงานต่อหนึ่งหน่วยเวลาหรือกำลัง $P = W/t$ ในระบบชีวภาพเกี่ยวข้องโดยตรงกับอัตราเมแทบอลิซึมพื้นฐาน (Basal Metabolic Rate: BMR) ในขณะที่งานด้านเกษตรกรรมและวิทยาศาสตร์สิ่งแวดล้อม เครื่องกลอย่างง่าย เช่น รอก พื้นเอียง สกรูลำเลียงเมล็ดพันธุ์ และระบบปั๊มน้ำพลังงานแสงอาทิตย์ ล้วนได้รับการออกแบบโดยคำนึงถึงการได้เปรียบเชิงกลจริง (AMA) และประสิทธิภาพเชิงกล ($e = W_{\text{out}}/W_{\text{in}} \times 100\%$) เพื่อประหยัดพลังงานต้นทุน
\end{bioappbox}

\begin{labbox}
1. ตรวจสอบตำแหน่งศูนย์ (Zero Error) ของเครื่องมือวัดทุกชนิดก่อนเริ่มบันทึกข้อมูล
2. ทำการวัดซ้ำอย่างน้อย 3 ถึง 5 ครั้งในแต่ละจุดทดลองเพื่อคำนวณหาค่าเฉลี่ยและค่าเบี่ยงเบนมาตรฐาน
3. บันทึกผลการวัดพร้อมระบุหน่วยเอสไอ (SI Units) และค่าความคลาดเคลื่อนของการวัดเสมอ
\end{labbox}

\section{ตัวอย่างโจทย์และการวิเคราะห์เชิงฟิสิกส์}

\begin{examplebox}
3. พื้นเอียง (inclined plane) เป็นพื้นผิวราบที่เอียงทำมุมกับแนวราบ ใช้สำหรับเคลื่อนย้ายวัตถุขึ้นสู่ที่สูงโดยใช้แรงน้อยกว่าการยกขึ้นตรงๆ แต่ต้องเคลื่อนที่เป็นระยะทางไกลกว่า ตัวอย่างที่พบได้ทั่วไป เช่น ทางลาดสำหรับเข็นสัมภาระขึ้นรถบรรทุก ทางลาดสำหรับรถเข็นวีลแชร์ และถนนบนภูเขาที่ตัดเป็นทางลาดเพื่อลดความชัน

4. ล้อและเพลา (wheel and axle) ประกอบด้วยล้อทรงกลมขนาดใหญ่ที่ยึดติดแน่นกับเพลาทรงกระบอกขนาดเล็กที่อยู่ตรงกลาง เมื่อออกแรงหมุนที่ล้อ (รัศมีใหญ่) จะเกิดแรงบิดที่เพลา (รัศมีเล็ก) ทำให้ได้เปรียบเชิงกลในการหมุน ตัวอย่างเช่น พวงมาลัยรถยนต์ ลูกบิดประตู และแกนม้วนสายเบ็ดตกปลา

5. สกรู (screw) คือพื้นเอียงที่พันเป็นเกลียวรอบแกนทรงกระบอก การหมุนสกรูรอบแกนหนึ่งรอบเทียบเท่ากับการเคลื่อนที่ขึ้นไปตามพื้นเอียงเป็นระยะทางเท่ากับเส้นรอบวงของสกรู ในขณะที่วัตถุเคลื่อนที่ในแนวแกนเพียงระยะเท่ากับระยะพิตช์ (pitch) ของเกลียว สกรูจึงให้ค่าได้เปรียบเชิงกลสูงมาก ตัวอย่างเช่น น็อตสกรูยึดชิ้นส่วน แม่แรงยกรถยนต์ และใบพัดของสว่าน

6. ลิ่ม (wedge) คือเครื่องกลรูปทรงสามเหลี่ยมที่มีลักษณะคล้ายพื้นเอียงสองด้านประกบกัน ใช้สำหรับแยกหรือผ่าวัตถุออกจากกันโดยเปลี่ยนแรงกดในแนวหนึ่งให้เป็นแรงแยกในแนวตั้งฉาก ตัวอย่างเช่น ขวานผ่าฟืน มีด และหัวตะปู
\end{examplebox}

\section{แบบฝึกหัดทบทวนและคำถามท้ายบท}

\begin{enumerate}
    \item จงอธิบายความแตกต่างระหว่างงานและกำลัง พร้อมยกตัวอย่างประกอบ
    \item คนงานคนหนึ่งยกกระสอบข้าวมวล 40 กิโลกรัม ขึ้นสูง 2 เมตร ในเวลา 4 วินาที จงหากำลังที่ใช้
    \item ปั๊มน้ำตัวหนึ่งใช้เวลา 10 วินาที ในการสูบน้ำมวล 200 กิโลกรัม ขึ้นไปสูง 5 เมตร จงหากำลังของปั๊มน้ำในหน่วยวัตต์
    \item เครื่องยนต์มีกำลัง 30,000 วัตต์ ขับเคลื่อนรถให้แล่นด้วยความเร็วคงที่ 20 เมตรต่อวินาที จงหาแรงขับเคลื่อนของเครื่องยนต์
    \item จงแปลงกำลัง 100 แรงม้า ให้อยู่ในหน่วยวัตต์และกิโลวัตต์
    \item จงแปลงกำลัง 5,000 วัตต์ ให้อยู่ในหน่วยแรงม้า
\end{enumerate}

% สิ้นสุดบทเรียน
